\documentclass{aa}

\usepackage[colorlinks=true, allcolors=blue]{hyperref}
\usepackage{graphicx} % for \includegraphics
\usepackage[subrefformat=parens]{subcaption} % for subfigures
\usepackage[varg]{txfonts} % for A&A fonts
\usepackage{mathrsfs} % for \mathscr
\usepackage{cleveref} % for \cref
\usepackage{csquotes} % for \enquote

% remove journal name+copyright for preprint (delete for journal submission)
\makeatletter % interpret @ in commands
\renewcommand*{\aa@journalname}{} % see aa.cls
\renewcommand*{\aa@manuscriptname}{} % see aa.cls
\fancyhead[CE]{\aa@headfont \aa@headings} % same as odd pages from aa.cls
\makeatother

\newcommand\taurec{\tau_\text{rec}}
\newcommand\kmin{{k_\text{min}}}
\newcommand\kmax{{k_\text{max}}}
\newcommand\EB{Einstein–Boltzmann}
\newcommand\boldu{{\bf u}}
\newcommand\scrH{\mathscr{H}}
\newcommand\innerproduct[2]{\left\langle #1, #2 \right\rangle}
\newcommand\diff[1]{\mathrm{d}{#1}}
\newcommand\TODO[1]{\textcolor{red}{TODO: #1}}

\begin{document}

%\title{Faster \EB{} codes with Chebyshev interpolation}
\title{Spectral interpolation in \EB{} codes}
\authorrunning{Sletmoen, H.}
%\titlerunning{Faster \EB{} codes with Chebyshev interpolation}
\author{Herman Sletmoen\thanks{Email: \texttt{herman.sletmoen@astro.uio.no}}}
\institute{Institute of Theoretical Astrophysics, University of Oslo, PO Box 1029 Blindern, 0315 Oslo, Norway}
\date{Received XX / Accepted XX}
\abstract{
	\EB{} codes are fundamental for computing theoretical predictions from cosmological models, and rely heavily on interpolation in its independent variables:
	conformal time $\tau$, wavenumbers $k$ and multipoles $\ell$.
	We show that interpolating in $k$ and $\ell$ with Chebyshev polynomials instead of cubic splines improves performance and precision,
	so fewer perturbations and line-of-sight integrals can be solved explicitly higher precision.
	Chebyshev interpolation converges exponentially fast for smooth functions,
	which offers a unique way for approximation-free \EB{} codes to regain performance without complicating the code.
	By considering source functions $S(\tau, k)$ and angular power spectra $C_\ell$ for matter and cosmic microwave background temperature, polarization and lensing,
	we find $1$-$2$ orders of magnitude lower worst-case interpolation error for the same number of points,
	but the result depends heavily on the target function and desired accuracy level.
	\TODO{speedup?}
	%The precise improvement depends largely on the target function and desired precision level.
	%\TODO{Some numbers? Equivalently, consider fixed precision and say speedup?}
	%Tradeoffs compared to cubic splines include less choice in the placement of interpolation points and dependence on smoothness.
	Our results suggest that Chebyshev interpolation should be used more frequently in physical modeling of smooth functions.
	Chebyshev interpolation is implemented in SymBoltz (version 1.6), which is publicly available at \url{https://github.com/hersle/SymBoltz.jl}.
	\iffalse
	\TODO{old; update and rewrite}
    We show how \EB{} codes can be sped up by interpolating source functions $S(\tau, k)$ in $k$ using Chebyshev polynomials.
    This gives $S(\tau, k)$ to higher precision by solving fewer ordinary differential equations (ODEs) for the perturbation $k$-modes, which is the most expensive part of the computation.
    Compared to the traditional method of interpolating with cubic splines, Chebyshev interpolation typically yields sources accurate to 1-2 more digits for the same number of $k$-modes, although the exact difference depends on the particular source and $k$-grid used for cubic splines.
    We illustrate this with source functions for cosmic microwave background (CMB) temperature, polarization and lensing, and for the source function of the matter power spectrum.
	\TODO{speedup for CMB?}
    Chebyshev interpolation converges rapidly for smooth functions and is a particularly good fit to approximation-free \EB{} codes like SymBoltz that benefit more from solving fewer perturbation $k$-modes, whereas approximation-based codes like CLASS violate the smoothness in $k$ but reduce the speed of solving each perturbation mode with approximations.
    Coincidentally, the clustering of Chebyshev nodes near the endpoints aligns well with the denser sampling at small $k$ needed for most source functions.
	As the position of the Chebyshev nodes is fixed, it also leads to a simple set of precision parameters $(\kmin, \kmax, N_k)$, in contrast to the infinite freedom in choosing the $k$-grid for splines that typically lead to a more complex hyperparameters.
	The Chebyshev interpolation has been implemented in SymBoltz (version 1.6), which is publicly available at \url{https://github.com/hersle/SymBoltz.jl}.
	\fi
}

\keywords{methods: numerical – cosmology: theory}

\maketitle

\nolinenumbers

\section{Introduction}
\label{sec:intro}

\EB{} codes, such as SymBoltz \citep{sletmoenSymBoltzjlSymbolicnumericApproximationfree2026a}, CLASS \citep{blasCosmicLinearAnisotropy2011b} and CAMB \citep{lewisEfficientComputationCMB2000b},
are fundamental tools for computing theoretical predictions from cosmological models.
They solve large and stiff ordinary differential equations (ODEs) and many line-of-sight integrals for the background and perturbation equations.
This yields source functions $S(\tau)$ and $S(\tau, k)$ and transfer functions $\Delta_\ell(k)$
of conformal time $\tau$, perturbation wavenumber $k$ and angular multipole $\ell$.

Interpolation is one of the most effective optimizations that accelerate these codes.
Instead of solving ODEs and integrals for every $k$ and $\ell$,
they do this only on coarse grids of $k$ and $\ell$ and then interpolate to more values.
This works because the output functions are smooth, but it introduces interpolation errors.
The default options of the codes are tuned for a balance between performance and precision,
but they slow down significantly with finer precision parameters that use finer interpolation grids.
Since the pioneering COSMICS code \citep{maCosmologicalPerturbationTheory1995},
every code has interpolated with cubic splines.

While cubic splines are solid general-purpose interpolation methods,
they are outperformed by spectral methods such as interpolation with Chebyshev polynomials \citep{trefethenApproximationTheoryApproximation2019}.
Cubic splines interpolate a function through $N$ points with $n=N-1$ piecewise cubic polynomials with an interpolation error that falls algebraically as $n^{-4}$.
In contrast, Chebyshev interpolation fits a single degree-$n$ polynomial through all points with an error that can decay exponentially as $e^{-n}$.
This makes Chebyshev polynomials converges much Faster and able to interpolate to higher precision with fewer points.
It is also simple to implement with the barycentric interpolation formula \citep{berrutBarycentricLagrangeInterpolation2004},
and unlike equispaced polynomial interpolation it always converges \citep{trefethenSixMythsPolynomial}.
%For example, Chebyshev approximation has popularized and shown to be very powerful by the Chebfun package in MATLAB \citep{battlesExtensionMATLABContinuous2004}.
However, the remarkably fast convergence rate depends on the underlying function being smooth (infinitely differentiable with bounded derivatives),
which makes them very interesting in physics where this property often holds.

Traditional \EB{} codes such as CAMB and CLASS use approximation schemes to switch between different equations at different $\tau$, $k$ and $\ell$ \citep[particularly Figure 10]{blasCosmicLinearAnisotropy2011b}.
This speeds up their solution, but significantly complicates both the physical and numerical implementation.
Moreover, it breaks smoothness by introducing discontinuities or kinks at the switching points.
Cubic splines are well suited for interpolation in this case because they use piecewise polynomials and behave locally,
so a change somewhere in the interpolant has little impact on it elsewhere.
Chebyshev interpolation instead uses one global polynomial, and a fault in one point can slow the interpolation convergence over the entire domain.

Recently, however, approximation-free \EB{} such as SymBoltz \citep{sletmoenSymBoltzjlSymbolicnumericApproximationfree2026a}, DISCO-EB \citep{hahnDISCODJDifferentiableEinsteinBoltzmann2024} and PyCosmo \citep{refregierPyCosmoIntegratedCosmological2018} have emerged.
They integrate the same equations for all $\tau$, $k$ and $\ell$ and spend more effort per point, but obtain fully smooth solutions.
In this paper, we investigate whether approximation-free can leverage this smoothness together with the rapid convergence of Chebyshev interpolation to improve performance and precision.
We implement this feature for flat universes in SymBoltz\footnote{\url{https://hersle.github.io/SymBoltz.jl/}}.

\section{Computations in \EB{} codes}
\label{sec:sources}

The equations in \EB{} codes are solved in stages for $\tau$, $k$ and $\ell$,
each interpolating solutions of earlier stages.
Background functions $S(\tau)$ are interpolated in $\tau$;
perturbation sources $S(\tau, k)$ in $\tau$ and $k$;
and angular spectra $C_\ell$ in $\ell$.

First, background and thermodynamics equations are solved for variables $S(\tau)$ that depend only on $\tau$.
For example, the Friedmann, continuity and Peebles equations are solved for the Hubble function $\scrH(\tau)$, free electron fraction $X_e(\tau)$, optical depth $\kappa(\tau)$, visibility function $v(\tau)$ and lookback time $\chi(\tau) = \tau_0 - \tau$.

Second, perturbation equations are solved for sources $S(\tau, k)$ by integrating large ODEs in $\tau$ for each independent $k$.
For example, this provides sources for CMB temperature ($T$), polarization ($E$), lensing ($\psi$) and matter overdensity ($m$):
\begin{subequations}
	\begin{align}
		S^T    & = v \, \Bigg(\frac{\delta_\gamma}{4} + \Psi + \frac{\Pi_\gamma}{16}\Bigg) + e^{-\kappa} \, (\Psi + \Phi)^\prime + (v u_b)^\prime + \frac{3\big(v \Pi_\gamma\big)^{\prime\prime}}{16 k^2}, \\
		S^E    & = \frac{3 v \Pi_\gamma}{16 \, (k \chi)^2}, \\
		S^\psi & = (\Psi+\Phi) \, \frac{\chi_\text{rec}}{(\chi-\chi_\text{rec}) \chi} \, \theta(\tau - \taurec), \\
		S^m    & = \delta_m + 3 (1+w_m) \, \frac{\scrH}{k^2} \, \theta_m .
	\end{align}
	\label{eq:sources}%
\end{subequations}
\Cref{fig:sources} plots examples of these functions.
This requires the background solution to be interpolated in $\tau$.

Third, line-of-sight integrals for transfer function multipoles%
\begin{equation}
    \Delta^A_\ell(k) = \int_{\tau_i}^{\tau_0} \diff{\tau} \, S^A(\tau, k) j_\ell(k \chi)
    \label{eq:los}
\end{equation}
are evaluated for a given source $S^A(\tau, k)$,
which must be interpolated in $\tau$ and $k$.
In turn, this is used to compute angular spectra%
\begin{equation}
    C_\ell^{AB} \propto \int_\kmin^\kmax \diff{k} \, k^2 P_0(k) \Delta^A_\ell(k) \Delta^B_\ell(k) \quad \text{(up to an $\ell$-scaling)},
    \label{eq:spectra}
\end{equation}
which correlates two sources $S^A(\tau, k)$ and $S^B(\tau, k)$.

Most sources $S(\tau, k)$ are smooth in $k$, but $j_\ell(k\chi)$ makes $\Delta_\ell(k)$ oscillate rapidly in $k$,
so the integrand of $C_\ell$ must be sampled at many $k$.
A standard strategy is therefore to solve the expensive perturbations explicitly for only $10^2$-$10^3$ $k$-values,
and then interpolate $S(\tau, k)$ to $10^3$-$10^4$ $k$-values before line-of-sight integration.
This is $10\times$-$100\times$ faster than solving the perturbations explicitly for all $k$, as this typically dominates the runtime.

Most $C_\ell$ are also smooth in $\ell$,
so another standard trick is to line-of-sight integrate for some $\ell$,
and finally interpolate to every integer $\ell$.
This also leads to a significant speedup.

%Practically, interpolation in $k$ amounts to selecting a $k$-grid to solve the perturbations for,
%before interpolating to a finer $k$-grid for several fixed $\tau$-slices.
%Ideally, the same $k$-grid a grid of $k$ and $\ell$ on which to solve perturbations and line-of-sight integrals.
%Ideally want to use same $k$ and $\ell$ for all sources.

\begin{figure*}
	\centering
	\begin{subfigure}{0.24\textwidth}
		\includegraphics[width=\textwidth]{figures/ST.pdf}
		\subcaption{\label{fig:sourceT}Temperature}
	\end{subfigure}
	\hfill
	\begin{subfigure}{0.24\textwidth}
		\includegraphics[width=\textwidth]{figures/SE.pdf}
		\subcaption{Polarization}
	\end{subfigure}
	\hfill
	%\\\bigskip
	\begin{subfigure}{0.24\textwidth}
		\includegraphics[width=\textwidth]{figures/SL.pdf}
		\subcaption{Lensing}
	\end{subfigure}
	\hfill
	\begin{subfigure}{0.24\textwidth}
		\includegraphics[width=\textwidth]{figures/SM.pdf}
		\subcaption{Matter}
	\end{subfigure}
	\iffalse
	\includegraphics[width=\textwidth]{figures/S.pdf}
	\fi
	\caption{%
		\label{fig:sources}
		Examples of the source functions \eqref{eq:sources} for temperature, polarization, matter and lensing (up to some normalization).
		The two former are shown around recombination,
		and the two latter throughout cosmic time.
		%Examples of source functions for temperature and polarization around the time of recombination,
		%and for lensing and matter from the early to the late universe.
		%They are the sources $S(\tau, k)$ in \cref{eq:sources} up to some normalization.
		The vertical scale is linear in arbitrary units.
	}
\end{figure*}

\section{Chebyshev polynomial interpolation}
\label{sec:chebyshev}

%Similarly to how Fourier series have extremely nice properties for interpolation of smooth functions that are periodic,
%interpolation with 

%Polynomial interpolation fits a single degree-$n$ polynomial through $N = n+1$ interpolation points,
%unlike (cubic) spline interpolation that joins $n$ piecewise (cubic) polynomials.
%This is perfectly stable and has very attractive properties for approximating smooth functions \emph{if} one uses Chebyshev polynomials.
%In contrast, interpolation with equispaced polynomials is unstable.
%In addition, Chebyshev interpolation is remarkably simple and efficient with Fast Fourier Transforms (FFTs) and the barycentric interpolation formula.
%Similarly to how Fourier series interpolation is optimal for interpolation of smooth periodic functions,
%Chebyshev polynomial interpolation is the closest analogue for near-optimal interpolation of an arbitrary smooth on a finite domain.

%We briefly summarize polynomial interpolation from the simplest possible practical angle, using the barycentric interpolation formula.
%As is common in the literature, we assume the domain is $x \in [-1, 1]$; otherwise it is trivial to linearly transform the domain before/after interpolation.

Any function $f(x)$ on the domain $x \in [-1, 1]$%
\footnote{An arbitrary domain $y \in [a, b]$ is easily mapped to the canonical interpolation domain $x \in [-1, 1]$ with $x(y) = 2 ( y - (b+a)/2 ) / (b-a)$.}
%\footnote{It is customary to work with the canonical interpolation domain $x \in [-1, 1]$. An arbitrary domain $y \in [a, b]$ is easily mapped to this domain with the transformation $x(y) = 2 ( y - (b+a)/2 ) / (b-a)$ and its inverse.}
is interpolated by a polynomial $p_n(x)$ of order $n=N-1$ through $N$ points $(x_i, f_i)$:
\begin{equation}
	f(x) \approx p_n(x) = \sum_j f_j \, \ell_j(x)
	\quad \text{with} \quad
	\ell_j(x) = \frac{\prod_{k \neq j} \, (x-x_k)}{\prod_{k \neq j} \, (x_j-x_k)} .
\label{eq:interppoly}
\end{equation}
The Lagrange polynomials satisfy $\ell_j(x_k) = \delta_{jk}$, so $p_n(x_k) = f_k$.

\iffalse
By convention, we here consider all points to lie on the canonical domain $x \in [-1, 1]$.
It is trivial to map points from an arbitrary domain $y \in [a, b]$ with the affine transformation
\begin{equation}
	y(x) = \frac{b+a}{2} + \frac{b-a}{2} x
	\  \iff \ 
	x(y) = \frac{2}{b-a} \, \bigg( y - \frac{b+a}{2} \bigg).
\end{equation}
\fi

\TODO{error nodal polynomial $\rightarrow$ Chebyshev nodes}

A priori, the points $x_j$ can be placed freely.
However, it is well understood that if $x_j$ are uniformly spaced, then $p_n(x)$ diverges as $n \rightarrow \infty$ near the boundaries of the interpolation domain.
This is called Runge's phenomenon (demonstrated in \cref{fig:recequi}).
Fortunately, this is \emph{not} a general property of polynomial interpolation, but only of the (poor) placement of points.

The fix is to add more points near the ends.
In fact, it is guaranteed that $p_n(x) \rightarrow f(x)$
if the point density asymptotes to%
\begin{equation}
	\rho(x) \propto \frac{1}{\sqrt{1-x^2}}
	\quad \text{as} \quad
	n \rightarrow \infty.
\label{eq:density}
\end{equation}
A common such set of points is the (type-II) Chebyshev nodes:
\begin{equation}
	x_m = \cos \frac{m\pi}{n},
	\qquad
	m = 0,\, 1,\, \ldots\, n.
\label{eq:chebnodes2}
\end{equation}
As explained in \cref{sec:coeffs}, they are the roots of the type-I Chebyshev polynomials
and enable practical advantages such as Fast Fourier Transforms (FFTs) computations.
\TODO{finish}
%This clusters points near the boundaries $x = \pm 1$ to prevent divergence there, and distributes them uniformly near the center $x = 0$.

Furthermore, the $p_n(x) \rightarrow f(x)$ converges remarkably fast if $f(x)$ is a smooth function.
Then polynomial interpolation is very attractive with an interpolation error that drops exponentially as%
\begin{equation}
	\max_{-1 \leq x \leq 1} \left\lvert p_n(x) - f(x) \right\rvert \propto e^{-n}
	\quad \text{as} \quad
	n \rightarrow \infty.
\label{eq:convergence}
\end{equation}

The best way to evaluate $p_n(x)$ for interpolation is through the barycentric interpolation formula \citep{berrutBarycentricLagrangeInterpolation2004}:%
\begin{equation}
	p_n(x) = \sum_{j=0}^n \frac{w_j}{x-x_j} f_j \,\,\,\Bigg/\,\,\, \sum_{j=0}^n \frac{w_j}{x-x_j}.
\label{eq:bary}
\end{equation}
Here $w_j = 1 / l^\prime(x_j)$ are interpolation weights that depend only on the location of the points through $\ell(x) = \Pi_k (x - x_k)$.
They can be precomputed for a general scheme,
but for the Chebyshev nodes \eqref{eq:chebnodes2} one can instead use the even simpler formula
\begin{equation}
	w_j = \begin{cases} (-1)^j/2 & \text{if } j = 0 \text{ or } j = n, \\ (-1)^j & \text{otherwise}. \end{cases}
\end{equation}
Barycentric interpolation is both stable and takes only $O(n)$ operations per $x$.
Despite suspicious divisions by $x-x_j$, it is stable in floating point arithmetic even as $x \rightarrow x_j$
and needs only a simple exception to return $f_j$ if $x = x_j$ exactly \citep{highamNumericalStabilityBarycentric2004}.
In contrast, direct evaluation of the Lagrange form \eqref{eq:interppoly} is both unstable and requires $O(n^2)$ additions and multiplications per $x$.

It is also possible to expand $p_n(x) = \sum_k a_k T_k(x)$ in Chebyshev polynomials $T_n(x)$.
A Fast Fourier Transform gets the coefficients $a_k$ from the values $f_i$ in just $O(n \log n)$ operations.
This is useful to study the convergence of $a_k$ with $k$.
\TODO{appendix.}

In contrast, piecewise interpolation methods join many low-order polynomials between each interpolation point and converge more slowly.
Piecewise linear interpolation is the simplest type, with an error that drops as $n^{-2}$.
Cubic splines are a popular general-purpose interpolation method, with an error that falls as $n^{-4}$.
They can be problematic near the ends, where the spline assumes boundary conditions that does not always match the nature of the data points.
In return, splines work well for arbitrary point configurations, are somewhat faster to evaluate, and works better with localized imperfections in the data points.

Thus, Chebyshev interpolation is extremely attractive when:%
\begin{itemize}
\item $f(x)$ is known to be smooth, to gain the fast convergence \eqref{eq:convergence};
\item $f(x)$ can be queried at any $x$, to match the density \eqref{eq:density};
\item $f(x)$ is expensive to evaluate, to benefit from fewer samples.
\end{itemize}
These conditions are seldom met when working with observed data, for example.
But they are often satisfied in theoretical applications, and particularly in \EB{} codes,
where the (unapproximated) underlying physics is smooth and the code is free to solve for any input.
In the next sections, we systematically consider interpolation in each of the three independent variables:
conformal time $\tau$, wavenumbers $k$ and multipoles $\ell$.

\section{Interpolation in time}
\label{sec:time}

We do not use Chebyshev interpolation in $\tau$.
Nevertheless, we discuss this aspect to compare it to interpolation in $k$ and $\ell$.

The background and perturbation equations for each $k$-mode are ODEs $u^\prime(\tau) = f(u, \tau)$ in $\tau$.
Their solutions have inherently local features,
as the universe undergoes eras dominated by radiation, matter and dark energy,
and shorter events such as recombination (see \cref{fig:sourceT}) and reionization.
Adaptive ODE solvers automatically adjust their time step to capture such local features.

Moreover, ODE solvers feature specialized dense output methods for interpolating between (adaptive) time points.
This is usually custom-tailored to the solver in a way that maximizes accuracy and reuses calculations of $f(u, \tau)$ at intermediate steps taken by the solver.
SymBoltz defaults to the implicit Rodas5P ODE solver, which uses a specialized fourth-order interpolant \citep{steinebachConstructionRosenbrockWanner2023a}.
Even without specialized dense output,
the time derivative $u^\prime(\tau) = f(u, \tau)$ can be computed analytically \enquote{for free} from the solution $u(\tau)$ through the ODE definition.
One can then fall back cubic Hermite interpolation,
which takes both $u(\tau)$ and $u^\prime(\tau)$ into account for accurate interpolation.

The combination of local features, adaptive time stepping and dense output methods of ODE solvers makes Chebyshev interpolation less suitable for $\tau$-interpolation.
This also impacts performance less,
as it is cheap to save more time points when integrating an ODE.
To take use of Chebyshev polynomials in time seriously,
one could rather attempt to use spectral ODE methods.
They expand $u(\tau) = \sum_k a_k T_k(\tau)$ in Chebyshev polynomials $T_k(\tau)$,
and use the ODE equations to set up and solve a matrix equation for $a_k$.
This is outside the scope of this paper, however.%
%In contrast, interpolation in $k$ (for fixed $\tau$) and $\ell$ has a is more independent and global in nature,
%and is a much better fit for Chebyshev interpolation.

\section{Interpolation over wavenumbers}
\label{sec:wavenumbers}

Along $k$, however, Chebyshev polynomials are a good candidate for interpolation.
The perturbation ODEs for each $k$ are independent,
and contain no analytical $k$-derivatives for use with Hermite interpolation.
The solution also has a more global nature in $k$,
as the background affects all scales at a time $\tau$ to some extent.

To use Chebyshev $k$-interpolation, SymBoltz reads a wavenumber interval $[\kmin, \kmax]$ and a polynomial degree $n$.
It then scales the Chebyshev grid \eqref{eq:chebnodes2} to $[\kmin, \kmax]$ and solves the perturbation ODEs for every $k_j$ on the grid.
Finally, any source $S(\tau, k)$ is interpolated to any $k \in [\kmin, \kmax]$ by evaluating the barycentric formula \eqref{eq:bary} with $f_j = S(\tau, k_j)$ for some  $\tau$.

\begin{figure*}
	\centering
	\begin{subfigure}{0.32\textwidth}
		\includegraphics[width=\textwidth]{figures/recequi.pdf}
		\subcaption{Equispaced polynomial interpolation}
		\label{fig:recequi}
	\end{subfigure}
	\hfill
	\begin{subfigure}{0.32\textwidth}
		\includegraphics[width=\textwidth]{figures/reccheb.pdf}
		\subcaption{Chebyshev polynomial interpolation}
		\label{fig:reccheb}
	\end{subfigure}
	\hfill
	\begin{subfigure}{0.32\textwidth}
		\includegraphics[width=\textwidth]{figures/reccub.pdf}
		\subcaption{Cubic spline interpolation}
		\label{fig:reccubic}
	\end{subfigure}
	\caption{%
		\label{fig:rec}
		Temperature source function $S^T(\tau, k)$ in \cref{fig:sourceT} at recombination interpolated with \subref{fig:recequi} equispaced polynomials, \subref{fig:reccheb} Chebyshev polynomials and \subref{fig:reccubic} cubic splines with uniform spacing.
		The interpolated results are compared to explicit (non-interpolated) solutions of the perturbation ODEs for 500 $k$-modes.
	}
\end{figure*}

\Cref{fig:rec} shows the temperature source interpolated over $k$ at the recombination time $\tau = \taurec$ (peak of the visibility function).
It illustrates that interpolation with both Chebyshev polynomials and equispaced cubic splines are stable,
unlike equispaced polynomial interpolation (Runge's phenomenon). 
As $n$ increases, the error of Chebyshev interpolation falls quickly and fully converges to numerical noise with only 80 points.
Cubic splines interpolate with 3-4 orders of magnitude higher error with the same number of points,
and fail to reach the level of Chebyshev polynomials even with 400 points.
This shows the fast $e^{-n}$-convergence of Chebyshev interpolation for smooth functions,
compared to the algebraic $n^{-4}$-convergence for cubic splines.

\begin{figure*}
	\centering
	\iffalse
	\begin{subfigure}{0.41\textwidth}
		\includegraphics[width=\textwidth]{figures/heaterrTcheb.pdf}
		\subcaption{Temperature interpolated with Chebyshev polynomials}
	\end{subfigure}
	\hfill
	\begin{subfigure}{0.41\textwidth}
		\includegraphics[width=\textwidth]{figures/heaterrTcub.pdf}
		\subcaption{Temperature interpolated with cubic splines}
	\end{subfigure}
	\\\bigskip
	\begin{subfigure}{0.41\textwidth}
		\includegraphics[width=\textwidth]{figures/heaterrEcheb.pdf}
		\subcaption{Polarization interpolated with Chebyshev polynomials}
	\end{subfigure}
	\hfill
	\begin{subfigure}{0.41\textwidth}
		\includegraphics[width=\textwidth]{figures/heaterrEcub.pdf}
		\subcaption{Polarization interpolated with cubic splines}
	\end{subfigure}
	\\\bigskip
	\begin{subfigure}{0.41\textwidth}
		\includegraphics[width=\textwidth]{figures/heaterrMcheb.pdf}
		\subcaption{Matter interpolated with Chebyshev polynomials}
	\end{subfigure}
	\hfill
	\begin{subfigure}{0.41\textwidth}
		\includegraphics[width=\textwidth]{figures/heaterrMcub.pdf}
		\subcaption{Matter interpolated with cubic splines}
	\end{subfigure}
	\\\bigskip
	\begin{subfigure}{0.41\textwidth}
		\includegraphics[width=\textwidth]{figures/heaterrLcheb.pdf}
		\subcaption{Lensing interpolated with Chebyshev polynomials}
	\end{subfigure}
	\hfill
	\begin{subfigure}{0.41\textwidth}
		\includegraphics[width=\textwidth]{figures/heaterrLcub.pdf}
		\subcaption{Lensing interpolated with cubic splines}
	\end{subfigure}
	\fi
	\includegraphics[width=\textwidth]{figures/heaterr.pdf}
	\caption{%
		\label{fig:heaterrs}
		Interpolation error of temperature, polarization, matter and lensing source functions in \cref{fig:sources} (left to right)
		interpolated with Chebyshev polynomials, cubic splines through uniform points and cubic splines through Chebyshev points (top to bottom) from 100 $k$-points in $k \in [1, 2000]$.
		As in \cref{fig:rec}, the error is computed relative to 500 explicitly solved perturbation $k$-modes.
	}
\end{figure*}

\Cref{fig:heaterrs} generalizes the comparison between Chebyshev and cubic spline interpolation to several source functions $S(\tau, k)$ and the full $(\tau, k)$-space.
To rule out differences only due to the sampling points,
we also include cubic splines using the Chebyshev interpolation points.
For the temperature and polarization sources, Chebyshev interpolation is again several orders of magnitude more accurate around recombination at $a(\tau) \approx 10^{-3}$,
where they are most important to capture.
For the matter and lensing sources,
Chebyshev interpolation is also more accurate and distributes the error uniformly over $k$.
Cubic splines struggle particularly at the boundaries of the interpolation interval,
due to their boundary conditions not matching the underlying function.

\begin{figure*}
	\centering
	\iffalse
	\begin{subfigure}{0.49\textwidth}
		\includegraphics[width=\textwidth]{figures/convT.pdf}
		\subcaption{Temperature}
	\end{subfigure}
	\hfill
	\begin{subfigure}{0.49\textwidth}
		\includegraphics[width=\textwidth]{figures/convE.pdf}
		\subcaption{Polarization}
	\end{subfigure}
	\\\bigskip
	\begin{subfigure}{0.49\textwidth}
		\includegraphics[width=\textwidth]{figures/convM.pdf}
		\subcaption{Lensing}
	\end{subfigure}
	\hfill
	\begin{subfigure}{0.49\textwidth}
		\includegraphics[width=\textwidth]{figures/convL.pdf}
		\subcaption{Matter}
	\end{subfigure}
	\fi
	\includegraphics[width=\textwidth]{figures/conv.pdf}
	\caption{
		\label{fig:convergence}
		Convergence of the $L_2$ norm of the interpolation errors over $(\tau, k)$ in \cref{fig:heaterrs}
		for interpolation with Chebyshev polynomials, uniformly spaced cubic splines and cubic splines through the points used for Chebyshev interpolation.
	}
\end{figure*}

\Cref{fig:convergence} compresses the error over $(\tau, k)$ to one number with the $L_2$ norm,
and shows its convergence with more interpolation points.
Every error falls faster or much faster with Chebyshev interpolation,
again demonstrating exponential convergence.

\begin{figure*}
	\centering
	\iffalse
	\begin{subfigure}{0.49\textwidth}
		\includegraphics[width=\textwidth]{figures/coeffsT.pdf}
		\subcaption{Temperature}
		\label{fig:coeffsT}
	\end{subfigure}
	\hfill
	\begin{subfigure}{0.49\textwidth}
		\includegraphics[width=\textwidth]{figures/coeffsM.pdf}
		\subcaption{Matter}
		\label{fig:coeffsM}
	\end{subfigure}
	\fi
	\includegraphics[width=\textwidth]{figures/coeffs.pdf}
	\caption{%
		\label{fig:coeffs}
		Coefficients $a_m$ in the Chebyshev polynomial expansion $S(\tau, k) \approx \sum_{m=0}^{400} a_m(\tau) T_m(k)$ of the source functions in \cref{fig:sources}.
		The temperature and polarization sources are evaluated at recombination,
		and the matter and lensing sources today.
		The curves correspond to decreasing ODE solver tolerances for the perturbations,
		which lowers the noise floor for the solution.
	}
\end{figure*}

\Cref{fig:coeffs} studies the convergence of the Chebyshev coefficients $a_m$ in the expansion $p_n(x) = \sum_m a_m T_m(x)$ described in \cref{sec:coeffs}.
This is a useful diagnostic,
because convergence of $a_m$ implies convergence of $p_n(x)$,
and has no counterpart with splines.
%is not possible with spline interpolation, because the interpolation error is bounded by the neglected coefficients: $\max_x \lvert f(x) - p_n(x) \rvert \leq \sum_{k=n+1}^\infty \lvert a_k \rvert$.
The convergence undergoes three phases (here for $S^T$):

\begin{itemize}
\item
For $0 \lesssim n \lesssim 25$,
the interpolant is still resolving the features of the source function and $a_m$ has not yet begun to converge.

\item
For $25 \lesssim n \lesssim 75$,
the interpolant has captured the main features and the coefficients decay rapidly.
The steepest fall is a line in the log-linear plot,
reflecting $e^{-n}$-decay of the error.

\item
For $n \gtrsim 75$,
the interpolant has fully converged to the floor of numerical precision set by the ODE sovler tolerance,
and there is no reason to go to higher order.
This matches the convergence with around 80 points in \cref{fig:reccheb}.
\end{itemize}

Chebyshev interpolation in $k$ is thus highly appropriate.
It can interpolate with orders of magnitude lower error than cubic splines with the same number of points.
The improvement depends on the source and generally becomes more dramatic at higher precision levels,
where Chebyshev interpolation can reach the noise floor using many fewer points than cubic splines.

\section{Interpolation over multipoles}
\label{sec:multipoles}

Angular power spectra $C_\ell$ are also smooth functions of the multipole $\ell$.
Instead of computing every $C_\ell$ explicitly with line-of-sight integration for every $\ell$,
\EB{} codes do this only for some $\ell$ and then interpolate to all desired integer $\ell$.
This again makes Chebyshev interpolation a good candidate.

Unlike source functions that vary continuously with $k$,
interpolation in the discrete $\ell$ raises a conceptual nuance.
As $C_\ell$ is physically interpreted only for integer $\ell$,
other codes interpolate from an integer $\ell$-grid (e.g. using every low $\ell$ and every $30$th high $\ell$).
This works with cubic splines because their stability is better for arbitrary grid configurations.
%Because the $\ell$-spacing is in steps of 1, CLASS relies on recursion relations in $\ell$ to evaluate the (hyper)spherical Bessel functions.
But it is incompatible with Chebyshev interpolation, because the node placement \eqref{eq:chebnodes2} forces the interpolation grid to arbitrary $\ell$.
This means that the line-of-sight integrals \eqref{eq:los} and the spherical Bessel functions $j_\ell(x)$ must be evaluated for both integer and non-integer $\ell$.

Physically, the generalization to non-integer $\ell$ is unproblematic.
The spherical Bessel functions $j_\ell(x)$ are (one of the linearly independent) solutions to the ODE (for $\ell > 0$)
\begin{equation}
	x^2 \frac{\mathrm{d}^2 j_\ell}{\mathrm{d} x^2} + 2x \frac{\mathrm{d} j_\ell}{\mathrm{d} x} + \left( x^2 - \ell(\ell+1) \right) j_\ell = 0,
	\quad
	j_\ell(0) = 0.
\label{eq:besselode}
\end{equation}
This is well-defined and smooth in both $\ell$ and $x$, so interpolation from arbitrary to integer $\ell$ converges to the right result.
Other codes that use integer $\ell$ with cubic splines already rely on the same property implicitly.
Thus, the discreteness concern is only a conceptual subtlety,
and practically it only means that $j_\ell(x)$ must be evaluated with $\ell$ treated as a continuous variable.%
\footnote{Extending $j_\ell(x)$ to continuous $\ell$ is analogous to analytically continuing the factorial $n!$ to the gamma function $\Gamma(n+1)$, for example.}

SymBoltz (version 1.6) lifts the integer-$\ell$ restriction to allow any interpolation method.
Before line-of-sight integration, it precomputes $j_\ell(x)$ on the $\ell$-grid set by the interpolation scheme using the library Bessels.jl\footnote{\url{https://github.com/JuliaMath/Bessels.jl/}},
which computes $j_\ell(x)$ continuously in both $\ell$ and $x$.
This enables Chebyshev interpolation, and also cubic spline interpolation without the integer-$\ell$ limitation.

\begin{figure*}
	\centering
	\begin{subfigure}{0.32\textwidth}
		\includegraphics[width=\textwidth]{figures/cmbTT.pdf}
		\subcaption{Temperature ($C_\ell^\mathrm{TT}$)}
		\label{fig:cmbTT}
	\end{subfigure}
	\hfill
	\begin{subfigure}{0.32\textwidth}
		\includegraphics[width=\textwidth]{figures/cmbEE.pdf}
		\subcaption{Polarization ($C_\ell^\mathrm{EE}$)}
		\label{fig:cmbEE}
	\end{subfigure}
	\hfill
	\begin{subfigure}{0.32\textwidth}
		\includegraphics[width=\textwidth]{figures/cmbLL.pdf}
		\subcaption{Lensing ($C_\ell^\mathrm{\psi\psi}$)}
		\label{fig:cmbLL}
	\end{subfigure}
	\caption{%
		\label{fig:cmb}
		Angular temperature, polarization and lensing self-correlation spectra interpolated with cubic splines versus Chebyshev polynomials.
		The interpolation is compared to the result from explicit line-of-sight integration for every $\ell$.
		Interpolation happens in $(\ell, \ell^5 C_\ell)$-space, where we found all interpolation methods to behave well for all three spectra.
		The cubic spline uses a grid with 8, 9 and 55 points on $\ell \in [2, 10]$, $[10, 100]$ and $[100, 3000]$, respectively, with shared points at $\ell = 10$ and $\ell = 100$.
		Chebyshev interpolation uses a single degree-69 polynomial through the Chebyshev nodes \eqref{eq:chebnodes2}, rescaled to $\ell \in [2, 3000]$.
		Piecewise Chebyshev interpolation joins a degree-5 polynomial on $\ell \in [2, 10]$ with a degree-64 polynomial on $\ell \in [10, 3000]$, with a shared point at $\ell = 2$.
	}
\end{figure*}

\Cref{fig:cmb} compares $\ell$-interpolation performance of cubic splines versus Chebyshev polynomials.
To fix the computational cost, we restrain both methods to exactly 70 $\ell$-points.
For each method, we also require that the same $\ell$-grid interpolates all three spectra $C_\ell^\mathrm{TT}$, $C_\ell^\mathrm{EE}$ and $C_\ell^{\psi\psi}$, so the precomputed $j_\ell(x)$ is common.

We see that Chebyshev interpolation converges to 3-4 orders of magnitude smaller relative error for $l \gtrsim 100$ than cubic splines for $C_\ell^\mathrm{TT}$ and $C_\ell^\mathrm{EE}$.
For $C_\ell^\mathrm{\psi\psi}$, the cubic spline performs better than a single Chebyshev polynomial.
However, this is because the Limber approximation switches on for $\ell \geq 10$ and makes a kink in $C_\ell$ that disrupts the global polynomial interpolant.
Cubic splines resist this because their behavior is local, so the high-$\ell$ region is not disrupted.
To fix this, we make a piecewise Chebyshev interpolant with one polynomial for $\ell \in [2, 10]$ (Limber approximation off) and another independent polynomial for $\ell \in [10, 3000]$ (Limber approximation on).
This restores the performance of Chebyshev interpolation to the level of the cubic spline.
%which is already excellent for $C_\ell^{\mathrm{\psi\psi}}$ because it is much smoother than $C_\ell^{\mathrm{TT}}$ and $C_\ell^{\mathrm{EE}}$.

Chebyshev interpolation is also simpler.
Fixing the computational cost to exactly 70 $\ell$-points,
we first tried a cubic spline with a uniform $\ell$-grid,
but this led to large interpolation errors for small $\ell$.
After progressively refining the $\ell$-grid with trial and error,
we ended up with a high-medium-low density $\ell$-grid over the subintervals of $\ell \in [2, 10, 100, 3000]$.
In contrast, Chebyshev interpolation performed better with a straightforward $\ell$-grid.

In this section, we always used the Chebyshev interpolation in $k$ from \cref{sec:wavenumbers} in order to isolate the effect to $\ell$-interpolation.
This means the cubic spline $\ell$-interpolation also benefits from Chebyshev $k$-interpolation.
In \cref{sec:kl}, we show the effect of using each interpolation method in both $k$ and $\ell$.

Note that we have concentrated on a flat universe.
In curved and open universes,
the spherical Bessel functions generalize to hyperspherical Bessel functions \citep{lesgourguesFastAccurateCMB2014a,tramComputationHypersphericalBessel2017a}.
This only modifies the ODE \eqref{eq:besselode} with an extra curvature-dependent term,
so we believe it is possible to generalize our $\ell$-interpolation to non-flat universes.

We conclude that Chebyshev polynomials are excellent also for $\ell$-interpolation.
They reach up to several orders of magnitude lower error in typical CMB temperature compared to cubic splines with the same number of points.
It is also used in the emulator Capse.jl \citep{boniciCapsejlEfficientAutodifferentiable2024a} for $C_\ell(\theta)$ of cosmological parameters $\theta$.
However, there the main purpose is to train it on a smaller number of independent coefficients $a_n(\theta)$ in the Chebyshev expansion of $C_\ell(\theta) = \sum_n a_n(\theta) T_n(\ell)$,
instead of every $C_\ell(\theta)$.
%As a bonus, the coefficients $a_n(\theta)$ also accelerate the likelihood computation compared to $C_\ell(\theta)$ because it allows a significant part of the likelihood computation to be precomputed once.
This approach is closely related to ours, as one can convert between $a_n$ and $C_\ell$ at the Chebyshev $\ell$-nodes by Fast Fourier Transforms (FFTs).

\section{Interpolation over wavenumbers and multipoles}
\label{sec:kl}

\begin{figure*}
	\centering
	\includegraphics[width=\textwidth]{figures/kl.pdf}
	\caption{%
		\label{fig:kl}
		Root-mean-square relative error $\big(\sum_\ell (C_\ell^\text{TT}/\bar{C}_\ell^\text{TT}-1)^2/N\big)^\frac12$ of the CMB temperature power spectrum $C_\ell^\text{TT}$
		for $50 \leq \ell \leq 2500$,
		computed with either Chebyshev interpolation or cubic spline interpolation in both $k$ and $\ell$ using different numbers of points.
		The error is relative to a reference spectrum $\bar{C}_\ell^\text{TT}$ computed with Chebsyhev interpolation from 1025 $k$-points and 513 $\ell$-points.
	}
\end{figure*}

Finally, we consider the combined effect of interpolating in both $k$ and $\ell$ with Chebyshev polynomials for angular spectra $C_\ell$.

\Cref{fig:kl} shows that Chebyshev interpolation can deliver 1-2 orders of magnitude lower relative error in $C_\ell$.
Here we consider the CMB temperature power spectrum $C_\ell^\text{TT}$ for $50 \leq l \leq 2500$, where it is exceptionally smooth (see \cref{fig:cmbTT}).
To visualize the error as a function of the number of $k$-points and $\ell$-points,
we compute the root-mean-square of the relative error over $\ell$.
Chebyshev interpolation fully converges to the $10^{-4}$ level using only around 50 interpolation points in each variable,
while cubic splines do not go below the $10^{-3}$-$10^{-2}$-level even 3 times more points per variable.

\section{Extension strategies}
\label{sec:ext}

\TODO{domain transformation to redistribute points, how to sample with different densities}

\TODO{value transformation}

\TODO{piecewise}

\section{Conclusion}
\label{sec:conclusion}

\TODO{quantify speedup and/or precision improvement, e.g. for CMB TT spectrum with 1e-4 error?}

\TODO{compare to e.g. cosmicnet, emulators, neural networks, etc.?}

\TODO{investigate $u = 1/(k+1) \in [0, 1]$ combined with Chebyshev interpolation that does not use endpoints?}

We have shown that Chebyshev polynomials improves the interpolation performance over $k$ and $\ell$ in \EB{} codes.
Compared to cubic spline interpolation used in other codes,
Chebyshev polynomials interpolates with significantly lower error from an equal number of points.
Our implementation in SymBoltz solves the perturbations and line-of-sight integrals on Chebyshev grids in $k$ and $\ell$ (including non-integers),
and interpolates the results with the barycentric formula.

The numerical properties of Chebyshev interpolation pairs strongly with the physical properties of approximation-free \EB{} codes.
Here the smooth solutions combine with the exponential convergence of Chebyshev interpolation.
This gives approximation-free codes a unique way to gain speed without sacrificing simplicity.
Approximation-based codes instead switch equations in $\tau$ and $k$ to speed up the solution,
which breaks smoothness and complicates the physics and numerics.

%Approximation-based codes: speed up with approximations, less smooth and need more points at a great cost in modeling complexity.
%Approximation-free codes: more smooth and need more points, speed up by integrating fewer $k$-modes at no cost in modeling complexity.

If implemented carefully, it is possible that approximation-based codes also benefit from Chebyshev interpolation.
For example, we saw that the Limber approximation for $\ell \geq \ell_\text{Limber}$ slowed convergence of Chebyshev $\ell$-interpolation,
but restored it by interpolating piecewise for $\ell < \ell_\text{Limber}$ and $\ell \geq \ell_\text{Limber}$.
The same strategy could be used in $k$-space in approximation-based codes,
but could be problematic when the switching criterion varies with both $\tau$ and $k$.
This could slow the convergence of Chebyshev interpolation because it responds globally to changes in one point.
The error of cubic splines decays locally and faster in their vicinity and could be the better method in such cases,
but they will not attain the same convergence rate and simplicity as Chebyshev interpolation in approximation-free codes.

\TODO{}
There are several aspects of interpolation we have not considered here.
Adaptive.
Domain transformation in $x$ for stretching and redistributing points,
value transformation in $y$.

Interpolation is only one of several computational methods used in \EB{} codes.
Lower precision of methods such as line-of-sight quadrature and the Limber approximation could overshadow the high precision of Chebyshev interpolation.
Research into other high-accuracy methods is therefore necessary for the results to converge to the right physical result with respect to all precision parameters.
This includes line-of-sight integration beyond the trapezoidal method and lensing methods beyond the Limber approximation \TODO{beyond-Limber}.
We believe Chebyshev interpolation is one of the ingredients for designing \EB{} codes
with numerical algorithms that leverage the physical properties of the equations as much as possible.

\TODO{}
We have here focused on a flat universe.
We believe the results generalize to open and closed universes,
with generalized wavenumbers $q(k)$ and hyperspherical Bessel functions in $\ell$.

\iffalse
\begin{acknowledgements}
	I thank heltonmc for implementing arbitrary $\ell$ in Bessels.jl.
	I thank Steven G. Johnson for creating the FastChebInterp.jl package and for advocating its use to 
nterpolate smooth and expensive functions on the forum for the Julia programming language.%
	\footnote{\url{https://discourse.julialang.org/t/117640/11}}
\end{acknowledgements}
\fi

\bibliographystyle{aa}
\bibliography{paper}

\appendix
%\onecolumn

\section{Chebyshev coefficients}
\label{sec:coeffs}

Instead of interpolating with the barycentric formula in \cref{sec:chebyshev},
a function can be approximated with a series of Chebyshev polynomials.
This offers an alternative approach to Chebyshev interpolation that is particularly useful for error analysis.
We give a brief summary of this approach based on \cite{trefethenApproximationTheoryApproximation2019}.

\iffalse
\begin{figure}
	%\includegraphics[width=\linewidth]{figures/chebpoly.pdf}
    \includegraphics[width=\linewidth]{figures/chebpoints.pdf}
    \caption{%
        Chebyshev grids on $[\kmin, \kmax]$ with 2 to 80 points.
    }
\end{figure}
\fi

The Chebyshev polynomials of the first kind are defined as
\begin{equation}
	T_n(x) = \cos(n \arccos x)
	\quad \text{for} \quad
	%\theta(x) = \arccos x
	x \in [-1, 1].
\label{eq:chebpolydef}
\end{equation}
To see that $T_n(x)$ is a polynomial,
evaluate $T_0(x)$ and $T_1(x)$ explicitly and use the identities $\cos(n\theta\pm\theta) = \cos(n\theta) \cos(\theta) \mp \sin(n\theta) \sin(\theta)$ with $\theta = \arccos x$ to find the polynomial recurrence%
\begin{equation}
	T_0(x) = 1,
	\quad
	T_1(x) = x,
	\quad
	T_{n+1}(x) = 2 x T_n(x) - T_{n-1}(x).
\label{eq:recurrence}
\end{equation}
From the trigonometric definition, we see that $T_n(x) = 0$ when $n\theta = \frac{\pi}{2}(1+2m)$ for integer $m$, so its roots are the Chebyshev nodes of the first kind:
\begin{subequations}
\begin{equation}
	x_m = \cos \frac{\pi(1+2m)}{2n}
	\qquad (0 \leq m \leq n-1).
\label{eq:chebnodes1_app}
\end{equation}
Similarly, $\lvert T_n(x) \rvert = 1$ when $n\theta = m\pi$, so its extrema are the Chebyshev nodes of the second kind (including the endpoints $x = \pm 1$):
\begin{equation}
	x_m = \cos \frac{m\pi}{n}
	\qquad
	(0 \leq m \leq n).
\label{eq:chebnodes2_app}
\end{equation}
\end{subequations}
Both sets of points have the optimal density \eqref{eq:density} for robust polynomial interpolation.
As we will see soon, their $\cos$-dependence also opens the door to Fast Fourier Transform methods.

Chebyshev polynomials are an orthogonal basis for expanding a function $f(x)$ on $x \in [-1, 1]$ in a unique series%
\begin{equation}
	f(x) = \sum_{k=0}^\infty c_k T_k(x).
\label{eq:infseries}
\end{equation}
They are orthogonal with respect to the inner product
(change variables to $x = \cos \theta$ and use orthogonality of cosines)
\begin{equation}
	%\innerproduct{T_m(x)}{T_n(x)} =
	\int_{-1}^1 \frac{T_m(x) T_n(x)}{\sqrt{1-x^2}} \, \diff{x} =
	\begin{cases} 0 & \text{if $m \neq n$,} \\ \pi & \text{if $m=n=0$,} \\ \pi/2 & \text{if $m=n \geq 1$,} \end{cases}
\label{eq:innerprodcont}
\end{equation}
so the coefficients for representing a function $f(x)$ are
\begin{equation}
	%c_0 = \frac{1}{\pi} \int_{-1}^1 \frac{T_k(x) f(x)}{\sqrt{1-x^2}} \, \diff{x},
	%\qquad
	c_k = \frac{2}{\pi} \int_{-1}^1 \frac{T_k(x) f(x)}{\sqrt{1-x^2}} \, \diff{x},
\label{eq:coeffscont}
\end{equation}
and half of that for $c_0$.
These coefficients represent $f(x)$ in the infinite-dimensional space of \cref{eq:infseries}, and is not very practical.%

\subsection{Approximation with Chebyshev projection}

A more practical way to approximate $f(x)$ is to take the coefficients $c_k$,
but truncate the infinite series \eqref{eq:infseries} at a finite $n$:
\begin{equation}
	f(x) \approx \hat{p}_n(x) = \sum_{k=0}^n c_k T_k(x).
\label{eq:projection}
\end{equation}
This is called Chebyshev projection,
because dropping coefficients effectively projects $f(x)$ from an infinite-dimensional vector space down to a lower-dimensional polynomial space.
It requires $f(x)$ at arbitrary $x$ to compute the integrals \eqref{eq:coeffscont}.
%If the integrand is sampled at $n$ points, it amounts to $O(n^2)$ work to construct the approximation.
This method is sometimes appropriate when one has an expression for $f(x)$ and can evaluate the integrals analytically, for example.
However, it is impractical for interpolation purposes.

\subsection{Approximation with Chebyshev interpolation}

% https://www.buttenschoen.ca/MATH551/build/point-choice-4c36fbc36ea8e8d082339f77a23601f8.pdf
An even more practical method is to construct the polynomial $p_n(x)$ that passes through $f(x)$ at the Chebyshev nodes \eqref{eq:chebnodes2_app}.
Instead of using the barycentric formula in \cref{sec:chebyshev},
we can expand $p_n(x)$ in Chebyshev polynomials:
\begin{equation}
	f(x) \approx p_n(x) = \sum_{k=0}^n a_k T_k(x)
	\quad \text{such that} \quad
	p_n(x_m) = f(x_m).
\label{eq:interpolation}
\end{equation}
Analogously to the continuous inner product \eqref{eq:innerprodcont}, the Chebyshev polynomials also satisfy a discrete orthogonality relationship at the Chebyshev nodes \eqref{eq:chebnodes2_app}:
\begin{equation}
	%\innerproduct{T_i(x_k)}{T_j(x_k)} =
	\sideset{}{''}\sum_{k=0}^n T_i(x_k) T_j(x_k) =
	\begin{cases} 0 & \text{if $i \neq j$,} \\ n & \text{if $i = j \in \{0, n\}$,} \\ n/2 & \text{if $i = j \in [1,n-1]$,} \end{cases}
\label{eq:innerproddiscrete}
\end{equation}
where the double-primed sum means that the first and last terms are halved.
This means that the coefficients $a_k$ are picked out as
\begin{equation}
	a_k = \frac{2}{n} \sideset{}{''}\sum_{j=0}^n f(x_j) \cos \left( \frac{jk\pi}{n} \right) .
\label{eq:coeffsdiscrete}
\end{equation}
This is exactly a type-I Discrete Cosine Transform (DCT),
so the coefficients can be computed with only $O(n \log n)$ instead of $O(n^2)$ operations!
It is another reason for why Chebyshev points are particularly attractive for polynomial interpolation.
The DCT converts from the function values $f(x_j)$ of the barycentric approach to the coefficients $a_k$ of the series approach.

Once $a_k$ are known, the interpolating series \eqref{eq:interpolation} is easily computed for any $x$ with $O(n)$ operations.
A naive implementation is to iterate through every term while computing the recurrence \eqref{eq:recurrence} forwards, but this loses precision for $\lvert x \rvert \rightarrow 1$.
A stable and slightly faster way is to iterate backwards with the \cite{clenshawNoteSummationChebyshev1955} algorithm
with an intermediate variable $b_k(x)$:
\begin{subequations}
\begin{align}
	b_{n+1}(x) &= b_{n+2}(x) = 0, \\
	b_k(x) &= a_k + 2x b_{k+1}(x) - b_{k+2}(x) \quad \text{for $k = n, \ldots, 0$}, \\
	p_n(x) &= \textstyle{\frac12 \big(a_0 + b_0(x) - b_2(x)\big)}.
\end{align}
\end{subequations}

Interpolating with the Chebyshev series \eqref{eq:interpolation} instead of the barycentric formula \eqref{eq:bary}
is the method of choice when one wants the coefficients $a_k$.
They provide a useful diagnostic:
the neglected coefficients indicate the maximum interpolation error
\begin{equation}
	\max_{-1 \leq x \leq 1} \big\lvert f(x) - p_n(x) \big\rvert \leq \sum_{k=n+1}^\infty \lvert a_k \rvert,
\label{eq:coefferr}
\end{equation}
so convergence of $a_k$ implies convergence of the whole $p_n(x)$.

Chebyshev projection and Chebyshev interpolation are generally different methods with $\hat{p}_n(x) \neq p_n(x)$.
Their coefficients $c_k \neq a_k$ is related by aliasing formulas,
and they usually perform within a factor of 2 of each other \citep{trefethenApproximationTheoryApproximation2019}.
However, the latter is much more practical for interpolation from a finite set of points.
It is also possible to interpolate with the roots \eqref{eq:chebnodes1_app} instead of the extrema \eqref{eq:chebnodes2_app} \TODO{ref}.

\section{Implementation details}

\TODO{non-obvious implementation details:}
Pull requests%
\footnote{\url{https://github.com/hersle/SymBoltz.jl/pull/93}}%
\footnote{\url{https://github.com/hersle/SymBoltz.jl/pull/115}}

\TODO{}

To set up a Chebyshev interpolation object,
...
Also write about optional domain/value transformation.

To set up a piecewise Chebyshev interpolation object,
...

To evaluate the Barycentric interpolation formula \eqref{eq:bary} for any $x$,
we loop over $j$, add $q_j = w_j / (x-x_j)$ and $q_j f_j$ to two cumulative variables for the denominator and numerator, and finally perform the division.
In every iteration, we also check if $x = x_j$ exactly and immediately return $f_j$ instead if that is the case.
As explained in \cref{sec:chebyshev}, this is stable for any $x$, also as $x \rightarrow x_j$.

To obtain the Chebyshev coefficients,
...

To support both Chebyshev interpolation and cubic spline interpolation,
we use the multiple dispatch feature of Julia ...

\end{document}
